\section{La receta de datos: de los atributos de cada fuente al equilibrio de la cola larga}

Contar con el mining pipeline no basta para mezclar directamente todos los pares de oraciones en el entrenamiento. Las distintas fuentes de datos difieren enormemente en alineación humana, ruido, escala y dependencia de modelos; estos atributos determinan la correlación de los errores, la escalabilidad y el sesgo. La clase primero enumera los retos del modelo, después usa un diagrama del Transformer para establecer el contexto mínimo y, por último, coloca los cinco tipos de datos en una misma tabla comparativa.

\subsection{Tres tipos de retos en la etapa del modelo}

Esta sección continúa con el pipeline de datos y explica por qué tener más datos se convierte de inmediato en un problema de modelado. Al crecer el volumen de datos, el modelo debe manejar a la vez el aumento efectivo de datos, la interferencia entre idiomas y la ampliación de capacidad. Los idiomas de bajos recursos necesitan compartir representaciones con los de altos recursos para obtener transfer, pero compartir en exceso puede producir interference; aumentar la capacidad alivia la competencia, pero hace que las tareas escasas se sobreajusten más rápido.

\lecturefigure{slide-23-model-challenges.jpg}{Retos acoplados del aumento de datos, la interferencia entre idiomas y el tamaño del modelo}{Video de la clase de Stanford 00:23:36}

\begin{knowledgebox}{Lectura de la figura: los tres bullets forman una tensión, no son TODO independientes}
Data augmentation amplía la cobertura, pero puede introducir ruido generado por modelos; reducir la language interference suele depender de más capacidad especializada, pero eleva el riesgo de sobreajuste; ampliar el modelo, a su vez, aumenta la complejidad del entrenamiento y del enrutamiento. Los MoE, EOM y el curriculum que se ven más adelante son precisamente una respuesta conjunta a este triángulo de contradicciones.
\end{knowledgebox}

\subsection{La ruta de traducción del Transformer encoder-decoder}

NLLB usa un encoder-decoder translation model. El encoder codifica la oración fuente completa en una representación contextualizada, y el decoder predice la oración objetivo token por token, condicionado por la representación de la fuente y por el prefijo objetivo ya generado. A diferencia del prompting con decoder-only, esta estructura separa explícitamente la comprensión del lado fuente y la generación del lado objetivo.

\lecturefigure{slide-24-transformer-preliminaries.jpg}{Flujo encoder-decoder del Transformer en la traducción automática}{Video de la clase de Stanford 00:24:03}

\begin{knowledgebox}{Lectura de la figura: las pilas superior e inferior asumen condiciones distintas}
El encoder de arriba lee simultáneamente los tokens fuente y forma representaciones entre posiciones; el decoder de abajo usa una causal mask, solo puede ver el prefijo objetivo y accede al encoder mediante cross-attention. En el entrenamiento suele usarse teacher forcing, mientras que en la inferencia el modelo depende de los tokens que él mismo generó antes, por lo que los errores pueden acumularse paso a paso.
\end{knowledgebox}

Dada una secuencia fuente $x$ y una secuencia objetivo $y=(y_1,\ldots,y_T)$, el objetivo de generación condicional es
\[
\mathcal{L}_{\text{MT}}(x,y)=-\sum_{t=1}^{T}\log p_\theta(y_t\mid y_{<t},x).
\]
$\theta$ representa los parámetros del modelo, $y_{<t}$ el prefijo objetivo, $p_\theta$ la probabilidad condicional del siguiente token y $T$ la longitud del objetivo. Todas las diferencias entre fuentes de datos terminan modificando la distribución de los pares $(x,y)$ que aparecen en esta pérdida.

\subsection{Los cinco tipos de fuentes de datos no pueden compararse solo por cantidad}

La tabla de la clase va agregando, uno por uno, NLLB-Seed, public bitext, mined bitext, multilingual MT backtranslation y bilingual SMT backtranslation. La tabla final pide al lector comparar cuatro atributos: si hay alineación humana, si hay ruido, si la escala es limitada y si depende de modelos existentes.

\lecturefigure{slide-25-data-source-comparison.jpg}{Provenance, ruido, escala y dependencia de las fuentes de datos de entrenamiento de NLLB}{Video de la clase de Stanford 00:26:03}

\begin{knowledgebox}{Lectura de la figura: leer columna por columna importa más que memorizar los nombres fila por fila}
Los datos human-aligned suelen tener alta calidad pero escala limitada; el public bitext puede venir de diversas fuentes y tiene calidad desigual; los mined data son escalables pero dependen del encoder y del filtrado; la MMT/SMT backtranslation puede aprovechar datos monolingües, pero hereda el sesgo del teacher model. Ninguna fuente cumple a la vez con alineación humana, ausencia de ruido, escala ilimitada e independencia de modelos.
\end{knowledgebox}

\begin{knowledgebox}{Términos clave: fuentes de datos de entrenamiento}
\begin{tabularx}{\textwidth}{@{}L{0.19\textwidth}L{0.25\textwidth}X@{}}
\toprule
Fuente & Problema que resuelve & Riesgos principales y relación con el curso \\
\midrule
NLLB-Seed & Aporta pares de oraciones confiables para el arranque & Escala pequeña, centrado en el inglés, puede contener translationese \\
Public bitext & Reutiliza recursos bilingües existentes & Fuentes heterogéneas, con límites de limpieza y de licencia distintos \\
Mined bitext & Amplía el corpus paralelo a partir de páginas web & Depende de LID, del encoder y de umbrales; contiene ruido por pares mal emparejados \\
MMT backtranslation & Genera oraciones fuente sintéticas con un teacher multilingüe & El sesgo del teacher se comparte entre idiomas; calidad inestable en bajos recursos \\
SMT backtranslation & Genera oraciones fuente sintéticas con un modelo estadístico bilingüe & Puede aportar un sesgo inductivo distinto, pero su cobertura y calidad son limitadas \\
\bottomrule
\end{tabularx}
\end{knowledgebox}

\subsection{Distribución original: unos cuantos idiomas concentran la gran mayoría de los datos}

La tabla de fuentes anterior explicaba la naturaleza de los datos; esta sección examina además la distribución de cantidades. Al combinar solo los public y seed data, el volumen de entrenamiento entre idiomas abarca varios órdenes de magnitud. En la figura, los idiomas de altos recursos forman barras altas y los de bajos recursos quedan pegados a la base; si se muestrea directamente según la frecuencia original, las actualizaciones del modelo quedan casi siempre dominadas por los idiomas grandes.

\lecturefigure{slide-26-public-seed-data.jpg}{Los public y seed data presentan una distribución de cola larga extrema}{Video de la clase de Stanford 00:26:33}

\begin{knowledgebox}{Lectura de la figura: el eje vertical está en escala logarítmica; las barras bajas no van ligeramente atrás}
Un eje vertical logarítmico implica que una diferencia de altura entre barras vecinas puede corresponder a diez o cien veces más datos. Primero hay que comparar la extensión de la cola larga y después observar idiomas de bajos recursos concretos. Un muestreo uniforme simple reutiliza pocos ejemplos y provoca sobreajuste; muestrear según la frecuencia natural hace que casi no aparezcan, por lo que se requiere la intervención conjunta del aumento de datos y de un muestreo controlado.
\end{knowledgebox}

\begin{warningbox}{Equilibrar los datos no es copiar cada idioma hasta la misma cantidad}
El muestreo repetido no agrega información nueva; al contrario, hace que las direcciones de bajos recursos memoricen más rápido el conjunto de entrenamiento. La estrategia de equilibrio debe considerar a la vez la diversidad efectiva de los ejemplos, el ruido, la dirección, la duración del entrenamiento y la capacidad del modelo.
\end{warningbox}

\subsection{Mining y backtranslation amplían la cola larga}

A continuación se observa cómo la augmentation modifica la cola larga. Al incorporar mined bitext y backtranslation, muchas direcciones de bajos recursos obtienen más volumen de entrenamiento y la cola larga se eleva. Pero la figura sigue mostrando diferencias marcadas, lo que indica que la función de la augmentation es mejorar la cobertura, no colocar a todos los idiomas en las mismas condiciones de datos.

\lecturefigure{slide-27-mining-backtranslation.jpg}{Mining y backtranslation amplían los datos de entrenamiento de bajos recursos}{Video de la clase de Stanford 00:27:00}

\begin{knowledgebox}{Lectura de la figura: cuando una barra crece, también hay que preguntar de dónde vienen los datos}
Los primary data son corpus alineados por personas o públicos; los mined data provienen de oraciones similares encontradas en la web; los BT data los genera un modelo, que produce oraciones fuente sintéticas a partir de oraciones objetivo monolingües. Al apilar los tres colores el total aumenta, pero su confiabilidad y su estructura de errores son distintas. Si un idioma de bajos recursos depende sobre todo de BT, el modelo puede reforzar en un ciclo cerrado los errores fijos del teacher.
\end{knowledgebox}

La backtranslation parte de una oración monolingüe $y$ en el idioma objetivo, usa un modelo inverso $q_\phi(x\mid y)$ para generar una oración fuente sintética $\hat{x}$ y luego entrena el modelo directo:
\[
\hat{x}\sim q_\phi(x\mid y),\qquad
\mathcal{L}_{\text{BT}}=-\log p_\theta(y\mid\hat{x}).
\]
$q_\phi$ es el teacher/backward model, $p_\theta$ es el modelo directo que se va a entrenar y $\hat{x}$ es la synthetic source. Su valor está en aprovechar grandes cantidades de texto monolingüe en el idioma objetivo; su riesgo es que los errores de $q_\phi$ se conviertan en etiquetas de entrenamiento.

\begin{lstlisting}[caption={Flujo conceptual de backtranslation con provenance y control de calidad}]
for target_sentence in target_monolingual_corpus:
    synthetic_source = backward_model.translate(target_sentence)
    score = quality_filter(synthetic_source, target_sentence)
    if score >= language_specific_threshold:
        training_buffer.add(
            source=synthetic_source,
            target=target_sentence,
            provenance="backtranslation",
            teacher=backward_model.version,
        )
\end{lstlisting}

\teachervoice{El giro central de la clase es que, después de «meter todos» los tipos de datos, el problema no termina. Más datos traen ruido y una cola larga más complejos, y el modelo necesita manejar la interferencia entre idiomas y el sobreajuste en bajos recursos.}

\subsection{Resumen de la sección}

La receta de datos de NLLB se compone del seed humano, el bitext público, el mined bitext y dos tipos de backtranslation. Cada fuente intercambia de manera distinta calidad, escala y dependencia de modelos; mejorar la cola larga no equivale a equilibrar la información, y el sobreajuste debe seguir controlándose desde el modelo y la programación del entrenamiento.
